\documentclass[11pt]{article}

\usepackage[a4paper,margin=2.5cm]{geometry}
\usepackage[T1]{fontenc}
\usepackage[utf8]{inputenc}
\usepackage[ngerman]{babel}
\usepackage{lmodern}
\usepackage{setspace}

\setstretch{1.1}

\begin{document}

\begin{center}
    {\LARGE \textbf{Vertraulichkeitsvereinbarung}}\\[1cm]
\end{center}

\vspace{0.5cm}

zwischen

\vspace{0.5cm}

\textbf{PreciPoint GmbH}\\
Parkring 6\\
85748 Garching b. München\\

\vspace{0.3cm}

(im Folgenden „PreciPoint“)

\vspace{0.8cm}

und

\vspace{0.5cm}

\textbf{Rücker und Schindele Beratende Ingenieure GmbH}\\
München\\

\vspace{0.3cm}

(im Folgenden „Partner“)

\vspace{1cm}

\section*{Präambel}

Im Rahmen der Masterarbeit von Herrn \textbf{Moritz Diez} werden Produktinformationen, technische Daten sowie interne Systeminformationen der PreciPoint GmbH verwendet und fachlich mit dem betreuenden Partner Rücker und Schindele Beratende Ingenieure GmbH diskutiert.

Zweck dieser Vereinbarung ist der Schutz vertraulicher Informationen der PreciPoint GmbH, welche dem Partner im Rahmen der Betreuung und fachlichen Diskussion der Masterarbeit zugänglich gemacht werden.

\section*{1. Vertrauliche Informationen}

Als vertrauliche Informationen gelten sämtliche nicht öffentlich bekannten Informationen der PreciPoint GmbH, unabhängig von ihrer Form der Übermittlung.

Hierzu zählen insbesondere:

\begin{itemize}
    \item Produktinformationen,
    \item technische Dokumentationen,
    \item Systemarchitekturen,
    \item Anforderungen und Spezifikationen,
    \item Entwicklungsinformationen,
    \item interne Projektinformationen,
    \item Software-, Hardware- und Systemkonzepte,
    \item Screenshots, Modelle, Diagramme und Auswertungen.
\end{itemize}

\section*{2. Zweckbindung}

Die vertraulichen Informationen dürfen ausschließlich zur Betreuung, Diskussion und fachlichen Unterstützung der Masterarbeit von Herrn Moritz Diez verwendet werden.

Eine Nutzung zu anderen Zwecken sowie eine Weitergabe an Dritte ist ohne vorherige schriftliche Zustimmung der PreciPoint GmbH nicht zulässig.

\section*{3. Vertraulichkeitspflichten}

Der Partner verpflichtet sich,

\begin{itemize}
    \item sämtliche vertraulichen Informationen streng vertraulich zu behandeln,
    \item die Informationen ausschließlich denjenigen Personen zugänglich zu machen, die unmittelbar an der Betreuung der Masterarbeit beteiligt sind,
    \item angemessene technische und organisatorische Maßnahmen zum Schutz der Informationen zu treffen,
    \item vertrauliche Informationen weder ganz noch teilweise ohne Zustimmung von PreciPoint zu veröffentlichen oder weiterzugeben.
\end{itemize}

\section*{4. Ausnahmen}

Die Verpflichtungen dieser Vereinbarung gelten nicht für Informationen,

\begin{itemize}
    \item die zum Zeitpunkt der Offenlegung bereits öffentlich bekannt waren,
    \item die nachträglich ohne Verstoß gegen diese Vereinbarung öffentlich bekannt werden,
    \item die dem Partner bereits nachweislich bekannt waren,
    \item die rechtmäßig von Dritten erhalten wurden.
\end{itemize}

\section*{5. Laufzeit}

Diese Vereinbarung tritt mit Unterzeichnung in Kraft und gilt für einen Zeitraum von fünf (5) Jahren.

Die Verpflichtung zur vertraulichen Behandlung bereits offengelegter Informationen bleibt auch nach Abschluss der Masterarbeit bestehen.

\section*{6. Veröffentlichungen}

Veröffentlichungen oder Präsentationen, welche vertrauliche Informationen der PreciPoint GmbH enthalten, bedürfen der vorherigen Abstimmung mit PreciPoint.

\section*{7. Gerichtsstand}

Es gilt deutsches Recht. Gerichtsstand ist München.

\vspace{2cm}

\noindent
Garching b. München, \today

\vspace{2cm}

\begin{tabular}{p{7cm} p{7cm}}
\hrulefill & \hrulefill \\
PreciPoint GmbH & Rücker und Schindele Beratende Ingenieure GmbH \\
\end{tabular}

\vspace{2cm}

\begin{tabular}{p{7cm}}
\hrulefill \\
Moritz Diez \\
\end{tabular}

\end{document}